\documentclass[11pt,a4paper]{article}
\usepackage[top=18mm,bottom=19mm,left=22mm,right=22mm]{geometry}
\usepackage[T1]{fontenc}
\usepackage{lmodern,amsmath,booktabs,xcolor,hyperref,csvsimple,pgfplots}
\pgfplotsset{compat=1.18}
\definecolor{ink}{HTML}{20343D}
\definecolor{sieve}{HTML}{08766D}
\definecolor{shuffle}{HTML}{A94B17}
\hypersetup{colorlinks=true,urlcolor=sieve,
  pdftitle={When should a recursive sieve hand off to pattern search?}}
\setlength{\parindent}{0pt}
\setlength{\parskip}{4pt}
\renewcommand{\arraystretch}{1.12}
\newcommand{\heading}[1]{\par\medskip{\large\sffamily\bfseries\color{ink}#1}\par}
\newcommand{\subhead}[1]{\par\smallskip{\sffamily\bfseries\color{ink}#1}\par}

% Read the primary curves directly from the archived output. Columns 1–6 are
% pilot, layout, paid call, number of pairs, recursive hit rate, random hit rate.
\def\recsolid{}\def\recisolated{}\def\randsolid{}\def\randisolated{}
\csvreader{../../results/handoff/hit_curve.csv}
  {1=\pilotlength,2=\keylayout,3=\paidcall,5=\rechit,6=\randhit}{%
  \ifnum\pilotlength=16\relax
    \expandafter\xdef\csname rec\keylayout\endcsname{\csname rec\keylayout\endcsname(\paidcall,\rechit)}
    \expandafter\xdef\csname rand\keylayout\endcsname{\csname rand\keylayout\endcsname(\paidcall,\randhit)}
  \fi}
\newcommand{\panel}[2]{%
  \begin{tikzpicture}
  \begin{axis}[width=.46\linewidth,height=43mm,scale only axis=false,
    title={#2},title style={font=\small\sffamily\bfseries},
    xmin=1,xmax=65,ymin=0,ymax=1,xtick={1,16,32,48,65},
    ytick={0,.25,.5,.75,1},yticklabels={0,25,50,75,100},
    tick label style={font=\footnotesize},xlabel={Paid objective calls},
    xlabel style={font=\footnotesize},ylabel={Hit by call (\%)},ylabel style={font=\footnotesize},
    grid=major,grid style={gray!20},axis line style={gray!65},tick align=outside]
  \draw[gray,densely dotted] (axis cs:16.5,0) -- (axis cs:16.5,1);
  \edef\drawrec{\noexpand\addplot[const plot,no marks,sieve,line width=1.2pt]
    coordinates {\csname rec#1\endcsname};}\drawrec
  \edef\drawrand{\noexpand\addplot[const plot,no marks,shuffle,dashed,line width=1.05pt]
    coordinates {\csname rand#1\endcsname};}\drawrand
  \end{axis}\end{tikzpicture}}

\begin{document}
{\sffamily\bfseries\fontsize{22}{25}\selectfont\color{ink}\raggedright
When should a recursive sieve hand off to pattern search?\par}
{\large Clustered and isolated targets on matched QAOA energy grids\par}
{\small Research draft for Pauli Rikula \hfill 21 November 2026\par}
\smallskip{\color{sieve}\rule{\linewidth}{0.9pt}}

\heading{The starting order changes which targets are found}
We compared two ways to spend the first 16 calls of a 65-call search. One visited
keys in recursive dyadic order and the other drew keys in random order without
replacement. Both then ran the same local pattern search from their best point.
On one solid block of targets, the recursive start found a target by call 65 in
88.6\% of placements, against 73.1\% for random. Isolating the same number of
targets reversed the ranking (77.5\% against 81.6\%). Each comparison paired
2,304 placements across 24 evaluation graphs.

\subhead{Table 1. First hits on the evaluation graphs}
\begin{center}\small
\begin{tabular}{llrrrr}
\toprule
& & \multicolumn{2}{c}{Hit by call 65} & \multicolumn{2}{c}{Difference in mean loss (calls)}\\
\cmidrule(lr){3-4}\cmidrule(lr){5-6}
Pilot & Target layout & Recursive & Random & Rec.\ $-$ random & 95\% interval\\
\midrule
16 & Solid block     & 88.6\% & 73.1\% & $-9.91$ & $[-11.54,-8.40]$\\
16 & Isolated keys   & 77.5\% & 81.6\% & $+2.03$ & $[+0.84,+3.16]$\\
\addlinespace
20 & Solid block     & 91.9\% & 78.4\% & $-9.25$ & $[-10.53,-8.03]$\\
20 & Isolated keys   & 81.0\% & 84.3\% & $+1.49$ & $[+0.38,+2.58]$\\
\bottomrule
\end{tabular}
\end{center}
{\small Loss is the first target call, assigning 66 to a miss. Negative differences
favor the recursive start. The 16-call comparison was fixed in advance.
Calibration selected 20 on separate graphs. Intervals resample whole graphs.\par}

\subhead{Figure 1. Pace of discovery with a 16-call pilot}
\begin{center}
\panel{solid}{One solid block}\hfill\panel{isolated}{Isolated keys}\\[-1pt]
{\small\sffamily\color{sieve}\rule[2pt]{12pt}{1.3pt} Recursive start
\qquad\color{shuffle}\tikz{\draw[dashed,line width=1.05pt] (0,0)--(.45,0);} Random start}
\end{center}
{\small The dotted line separates the pilot from the first local poll on call 17.
Each curve includes every placement through call 65, including early-converged
misses. On solid blocks, most of the recursive advantage was already present
at handoff. On isolated keys, the gap widened during local polling.\par}

\newpage
\heading{Measuring the handoff}
We used two disjoint cohorts of 24 MaxCut graphs with 7--9 vertices. For each
graph, an exact one-layer QAOA simulation supplied 256 energies on a periodic
$16\times16$ angle grid. For each of three initial points, the target threshold
was $E_0-0.8(E_0-E_{\min})$ with tolerance $10^{-12}$. Each of 32 paired draws
kept the initial energy, the multiset of energies and the target count fixed.
The target energies occupied either one uniformly placed block or a uniform
set of pairwise nonconsecutive keys in linear Morton order. Key zero was the
common invalid start. Reindexing changes the angular neighborhoods, so the
comparison measures performance on these planted layouts.

Call 1 paid for the start. At handoff each method retained its evaluated
incumbent and polled four wrapped angle neighbors per step, halving the step
after an unsuccessful group. Repeated polls counted. At a 16-call recursive
handoff the four initial local polls revisited the coarse lattice. The runner
scored the first successful evaluation even if polling continued afterward.

\subhead{Table 2. Stage of the first hit}
\begin{center}\small
\begin{tabular}{llrrrrrr}
\toprule
&&\multicolumn{3}{c}{Recursive start}&\multicolumn{3}{c}{Random start}\\
\cmidrule(lr){3-5}\cmidrule(lr){6-8}
Pilot & Layout & Pilot & Local & Miss & Pilot & Local & Miss\\
\midrule
16 & Solid    & 2,005 & 36  & 263 & 1,417 & 267 & 620\\
16 & Isolated & 1,416 & 370 & 518 & 1,435 & 444 & 425\\
20 & Solid    & 2,085 & 33  & 186 & 1,610 & 196 & 498\\
20 & Isolated & 1,594 & 272 & 438 & 1,601 & 342 & 361\\
\bottomrule
\end{tabular}
\end{center}
At 16 calls, polling added 36 first hits on 299 solid cases still at risk after
the recursive pilot. With random initialization it added 267 on 887. On
isolated keys the corresponding counts were 370 on 888 and 444 on 869.
Table~1 intervals use 2,000 paired bootstrap resamples of graphs within the six
size/density strata, keeping starts and placements together.

\heading{How long to keep sieving}
When at least one consecutive target run has length $L$, the zero-first
recursive sieve on $N=2^n$ keys guarantees a hit by
$D(L)=N/2^{\lfloor\log_2 L\rfloor}$ calls. The first $2^j$ probes are spaced
$N/2^j$ apart, so every run of that length contains a probe.
All 1,760 evaluation blocks of width 16, 24 or 32 were hit in the first
16 recursive calls. On the 544 width-eight blocks the 16-call hybrid had
245 pilot hits, 36 local hits and 263 misses. Continuing the sieve through
call 32 would guarantee a hit on every solid placement in this evaluation set.
For isolated targets, $L=1$ gives a deadline of 256.

For unknown layout, calibration minimized the mean of $\min(T_k,66)$ over both
layouts. For a pilot of length $k$, $T_k$ counts calls through the first
target in either stage ($T_k=\infty$ on a miss). For $k=12,14,16,18,20$, the respective
mean losses were 23.57, 21.52, 19.35, 18.45 and 17.69. Twenty was the
best tested cutoff and the upper endpoint. The loss at longer cutoffs remains open.

{\footnotesize\smallskip
\textbf{Reproduction.} The \href{https://github.com/kummahiih/variational-quantum-algorithms/tree/cf6758371ddac34910474ade050eb5357326376f/results/handoff}{archived run records and per-call curves}
and \href{https://github.com/kummahiih/variational-quantum-algorithms/blob/8c4acad88fc133a4b9d4ece7675024ae7a4ad767/paper/handoff/README.md}{prespecified comparison and runner instructions} are public.
The \href{https://github.com/kummahiih/variational-quantum-algorithms/tree/a3444d4b071942ead4f008fb10b61a61e7449a4c/paper/clusters}{companion study} describes the earlier 13-call comparison.
The project owner supplied the simulation output. OpenAI Codex assisted with
implementation, analysis and manuscript drafting. This draft awaits author review.\par}
\end{document}
